% !TEX root = ../algebra_02.tex

\section{Поле разложения многочлена. Существование}

\begin{Definition}{Поле разложения}{}
	Пусть $K$ --- поле, $f \in K[x]$ --- ненулевой многочлен. Расширение $L/K$ называется \textbf{полем разложения} многочлена $f$, если выполняются два условия:
	\begin{enumerate}
		\item В кольце многочленов $L[x]$ многочлен $f$ раскладывается на линейные множители:
		\[ f(x) = c(x-a_1)\dots(x-a_n), \quad \text{где } a_i \in L \]
		\item Поле $L$ порождается над $K$ всеми корнями этого многочлена:
		\[ L = K(a_1, \dots, a_n) \]
	\end{enumerate}
\end{Definition}

\begin{Theorem}{Существование поля разложения}{}
	Пусть $f \in K[x]$. Тогда существует поле разложения многочлена $f$.
\end{Theorem}

\begin{proof}
	Докажем индукцией по $\deg f$.

	\textbf{База:} При $\deg f = 1$ в качестве поля разложения подходит само поле $L = K$.

	\textbf{Шаг:} Пусть $\deg f > 1$. Пусть $f_0 \mid f$ --- неприводимый делитель в кольце $K[x]$. Построим факторкольцо (которое является полем в силу неприводимости $f_0$):
	\[ K_1 := K[x]/(f_0) \]
	Обозначим класс $x$ в $K_1$ как $x_1 = \overline{x}$. Тогда по определению $f_0(x_1) = 0$, откуда следует, что и $f(x_1) = 0$.
	Значит, по теореме Безу в $K_1[x]$ многочлен $f$ представим в виде:
	\[ f = (x - x_1)g, \quad g \in K_1[x] \]
	По предположению индукции, примененному к многочлену $g \in K_1[x]$, существует поле разложения $L/K_1$. Тогда в $L[x]$ многочлен $g$ раскладывается на линейные множители:
	\[ g = \varepsilon(x - x_2) \dots (x - x_n) \]
	Следовательно, в $L[x]$ многочлен $f$ также полностью раскладывается:
	\[ f = \varepsilon(x - x_1)(x - x_2) \dots (x - x_n) \]
	Само поле разложения порождается всеми корнями:
	\[ L = K_1(x_2, \dots, x_n) = K(x_1)(x_2, \dots, x_n) = K(x_1, x_2, \dots, x_n) \]
\end{proof}
